% Checked continuation after checkpoint 02; included into paper_B.tex by script.
\subsubsection*{Statement B.1e (Proved): the regular target over an imperfect field}
Let $k$ be any field of characteristic $p>0$.  Let $f:X\to Y$ be a
proper birational morphism of proper geometrically integral $k$-schemes,
with $X$ geometrically normal and $Y$ regular.  Then pullback is an
isomorphism on $\operatorname{Pic}^{\tau}$.
\begin{proof}
The equality $f_*\mathcal O_X=\mathcal O_Y$ commutes with field extension.
Since $X_{\bar k}$ is normal, this equality shows that $Y_{\bar k}$ is
normal as well: a rational function integral over an affine open of
$Y_{\bar k}$ pulls back to a regular function on its inverse image in
$X_{\bar k}$ and then descends by the displayed equality.  Thus both
Picard torsion components are proper and of finite type, and pullback
is a closed immersion $Q\hookrightarrow P$ as in B.1.

The smooth locus of $X/k$ is nonempty, so it has a point over some
finite separable extension $l/k$.  Normality of $X$ and regularity of
$Y$ survive that extension.  It suffices to prove the assertion over
$l$, and we therefore assume that $X(k)\ne\varnothing$; its image gives
a point of $Y$ as well.  For every finite commutative $k$-group $H$,
Raynaud's Proposition 6.2.1 and the sections of the structure maps give
\[
 \operatorname{Hom}(H,\operatorname{Pic}_{Z/k})
 \simeq H^1(Z,H^D)/H^1(k,H^D),\qquad Z=X,Y.
\]
Indeed, the Leray edge map is surjective because the section makes
$H^2(k,H^D)\to H^2(Z,H^D)$ injective.  B.1b therefore implies that
$\operatorname{Hom}(H,Q)\to\operatorname{Hom}(H,P)$ is bijective.

We can detect the closed subgroup over $k$ itself.  For each positive
integer $n$, the group $P[n]$ is finite: after an algebraic closure
this follows from finiteness of multiplication on the reduced abelian
identity component and the finitely many components.  Hence $P[n]$
lies in $Q$.  After extending to $\bar k$, the torsion-density and
torsion-representative argument of B.1c gives
$(P_{\bar k})_{\mathrm{red}}\subset Q_{\bar k}$.
The finite kernel of every relative Frobenius power of $P$ also lies
in $Q$.  Choose a power killing the nilradical of $P_{\bar k}$ and use
the final fppf argument of B.1c.  It gives $Q_{\bar k}=P_{\bar k}$,
hence $Q=P$.  We never apply normality of $X$ after a purely
inseparable extension in the torsor-purity argument.
\end{proof}

\subsubsection*{Statement B.2h (Proved): full stabilization over a perfect field}
Let $k$ be perfect of characteristic $p>0$, and let $K/k$ be finitely
generated.  The direct system of $\operatorname{Pic}^{\tau}_{X/k}$ over
all proper normal models of $K/k$ is essentially constant.
\begin{proof}
Let $k_0$ be the algebraic closure of $k$ in $K$.  It is a finite
separable extension.  Every proper normal model has its canonical
$k_0$-structure because its ring of global functions is $k_0$, and
every domination preserving the marking by $K$ is $k_0$-linear.
An isomorphism of the Picard schemes over $k_0$ gives an isomorphism
over $k$, as can be checked after extending to $\bar k$ and taking
the product over the embeddings of $k_0$.  We may therefore assume
that $k$ is algebraically closed in $K$.

By B.2g choose a normal projective model $S$ of $K\bar k/\bar k$
after which the geometric system is constant.  Its equations and
function-field marking descend to a finite extension of $k$.  Since
$k$ is perfect, enlarge this to a finite Galois extension $l/k$.
Write $S_l$ for the descended model and $\Delta=\operatorname{Gal}(l/k)$.
Take the closure of the rational diagonal with function field $Kl$
in the product of the $\Delta$-conjugates of $S_l$, and normalize it.
The resulting projective normal $l$-scheme $N_l$ has the canonical
semilinear $\Delta$-action and dominates every conjugate of $S_l$.
Its descent datum is effective by
\cite[Lemma 39.25.3, Tag 0CCJ]{StacksB}; equivalently, one can descend
a tensor product of the conjugate ample line bundles and its graded
section algebra.  We obtain a normal projective $k$-model $N$ of $K$.
Normality is geometric here because the ground fields are perfect.

For every normal proper $W\to N$, the base change $W_{\bar k}$ is
normal and dominates $S$.  Thus the map
$\operatorname{Pic}^{\tau}_{N/k}\to\operatorname{Pic}^{\tau}_{W/k}$
becomes an isomorphism over $\bar k$ and is therefore an isomorphism.
Normalized graphs give a common normal model of any two proper
models over the perfect field.  Consequently the models dominating
$N$ form a cofinal subsystem, and the isomorphisms just obtained
give essential constancy in the usual filtered sense.
\end{proof}

\subsubsection*{Statement B.2i (Refuted): filteredness over an imperfect field}
The category of marked proper geometrically normal models of a fixed
function field need not be filtered under domination.  This failure
already occurs for two smooth projective rational threefolds over an
imperfect field.  It invalidates a normalized-graph cofinality argument
in that generality; it does not itself refute a Picard-constancy
assertion formulated for a nonfiltered category.
\begin{proof}
Let $k=\mathbb F_p(a,b)$, with $a,b$ algebraically independent, and
let $K=k(x,y,z)$.  Take $X=\mathbb P^2_k\times\mathbb P^1_k$ with
its standard marking.  Take a second copy $Y$ with the marking given
on dense affine charts by
\[
 (x,y,z)\dashmapsto (x,y,w),\qquad
 w=z\frac{x^p-a}{y^p-b}.
\]
The inverse multiplies $w$ by $(y^p-b)/(x^p-a)$, so the map is
birational.  Both models are smooth and projective over $k$.
Put $f=x^p-a$ and $g=y^p-b$.  On the open where $z,w$ are invertible,
the closure of the graph has coordinate ring
\[
 R=B[w,w^{-1}],\qquad
 B=k[x,y,t,t^{-1}]/(y^p-b-t(x^p-a)),\quad t=z/w.
 \tag{B.2.5}
\]
The polynomial in $y$ is irreducible over $k(x,t)$, since
$b+t(x^p-a)$ is not a $p$-th power: its $t$-derivative is $x^p-a\ne0$.
Thus these rings are domains.

We verify that $B$, and hence $R$, is regular over $k$.  Where
$f\ne0$, the derivative of the defining equation with respect to
$t$ is nonzero, so the hypersurface is smooth.  On $f=0$, its
quotient is
\[
 B/(f)=l[t,t^{-1}],\qquad l=k(a^{1/p},b^{1/p}),\quad [l:k]=p^2.
\]
The element $f$ is a nonzerodivisor.  At every prime above $(f)$,
the local ring modulo $f$ is regular; lifts of its regular parameters,
together with $f$, generate the maximal ideal in the dimension
required for regularity.  This proves the assertion.  In particular
the normalization of the entire graph leaves this open unchanged.

After extending to $\bar k$ and putting
$u=x-a^{1/p}$ and $v=y-b^{1/p}$, the same open becomes
\[
 \bar k[u,v,t,t^{-1},w,w^{-1}]/(v^p-tu^p).
\]
At the codimension-one prime $(u,v)$, the local ring has dimension
one and its maximal ideal needs the two generators $u,v$; the
relation has degree at least two.  It is not regular, and therefore
is not normal.  Explicitly, its normalization adjoins $s=v/u$, with
$s^p=t$, and is $\bar k[u,s,s^{-1},w,w^{-1}]$.

Suppose a proper geometrically normal model $Z$ dominated both
marked models.  Its map to $X\times_kY$ factors through the graph
and then through the normalization of the graph, by normality.
This is a proper birational morphism.  It is an isomorphism over
an open containing the generic point of the divisor $f=0$ in
(B.2.5), because the target is normal.  Consequently $Z$ has the
same geometrically nonnormal local open near that divisor,
contradicting geometric normality.  No common such refinement exists.
The accompanying symbolic certificate verifies the graph equation,
the inverse rational map, and the normalization substitutions for
$p=2,3,5$; the proof applies to every prime.
\end{proof}
